\newpage
\ifdefined\useChapters
\chapter{Traído}
\else
\chapter{}
\fi

Eu sempre imaginei que minha irmã não saberia lidar muito bem quando soubesse de tudo que venho passando desde que acordei daquele coma — mas nunca, nem nos meus cenários mais pessimistas, pensei que seria daquele jeito. Ser levada por um psicopata, amarrada, esfaqueada, curada de um jeito que nenhum ser humano deveria presenciar, e depois teleportada para outro lugar…  
Era absurdo.  
Era cruel.  
E era a pior maneira possível de revelar a verdade.

Mas agora não adianta me perder no arrependimento.  
Eu não posso ficar aqui, parado, lamentando não poder contar com Lúcia.  
O mundo está correndo para um destino que eu já vi — e ele não espera ninguém.

Preciso agir.  
Preciso me preparar.  
Preciso aprender o máximo de habilidades possível e voltar para impedir que o futuro aconteça daquele jeito.

E para isso, preciso do Ítalo.

Ou… achava que precisava.  
Ele parecia ser meu companheiro mais importante até aqui.  
Alguém que eu finalmente conseguia confiar.  
Alguém que me via como igual — coisa rara pra mim.

Mas agora que eu conheço muito mais o meu potencial, consigo sentir a mentira no ar.  
Há coisas no Ítalo que eu não quis ver antes.  
Detalhes que meu cérebro ignorou porque eu precisava de um aliado.

Agora consigo perceber melhor.

Antes, meu teletransporte era impulsivo, instável.  
Eu ia parar em lugares aleatórios, tempos desconhecidos, quase sempre no susto.  
Mas agora, não.  
Agora que consigo me conectar mentalmente com as pessoas, tudo amplificou.  
Eu localizo intenções, medos, ambiguidades — e consigo chegar onde quero com precisão absurda.

Mesmo assim, eu sabia: Ítalo ficava desconfortável quando eu aparecia simplesmente do nada.  
Então decidi chegar mais longe, a vinte metros acima do prédio. As pessoas raramente olham para cima. E invisível, menos ainda.

A cidade me recebeu com fumaça, o cheiro de queimado impregnado nas paredes.  
Condomínio de prédios… queimados.  
O mesmo lugar onde Pedro e Larissa moravam.

Coincidência?  
Não existe coincidência quando se trata da minha vida.

Fiquei invisível e comecei a volitar devagar, atravessando paredes como se fossem cortinas de tecido leve.  
E foi então que senti.

Ítalo estava ali.  
Próximo.  
Mas sua pintura mental era… estranha.  
Amarela berrante em alguns pontos, vermelha profunda em outros, e depois um cinza pálido — o tipo de cor que só aparece quando alguém tenta esconder algo com vergonha.

Eu segui o rastro.

Quando finalmente atravessei a última parede, dei de cara com os dois.

Larissa.  
E Ítalo.

Lado a lado.

— Agora já estamos aqui na sua casa, o que você quer? — perguntou Ítalo, olhando para todos os lados com inquietação visível.

Larissa cruzou os braços.

— Você está com medo de alguma coisa? Por que está olhando assim pra todo lado?

— Você não faz ideia do que o garoto é capaz — respondeu Ítalo, apertando os lábios. — A gente não tem a menor chance de sequer neutralizá-lo.

Larissa riu, mas era um riso duro, nervoso.

— Não entendo esse seu medo todo. Eu já enfrentei ele algumas vezes. Lembro do medo estampado no rosto dele.

— Isso foi antes — rebateu Ítalo, a voz carregada de tensão. — Agora ele está muito mais poderoso.  
Da última vez que vi, me disse que precisava ficar sozinho para entender até onde podia chegar.  
E eu imagino que tenha conseguido.  
Já vi ele fazer coisas que é até difícil de explicar em voz alta.

Meu estômago virou.

Eu não acredito.  
O Ítalo…  
o único amigo que achei que tinha…  
estava me enganando esse tempo inteiro?

Me fez de tolo.  
Me distraiu.  
Enquanto isso… os três tinham outro plano.

Observei os dois tentando cochichar inútil e desesperadamente, como se eu não pudesse ouvir.  
Eu podia.  
E podia ver.

E agora tudo fazia sentido.

Da pintura mental que eu via em Ítalo — aquela mistura de medo, culpa e vergonha — até o brilho determinado nos pensamentos de Larissa.

Eles queriam juntar pessoas despertas.  
Gente de todas as idades, de várias raças, gêneros, religiões, convicções.  
Queriam fazer uma transmissão ao vivo, revelando tudo para o mundo.  
Mostrar que o poder não era privilégio de poucos.  
Mostrar que *todos* podiam despertar.

Um plano infantil, talvez.  
Mas corajoso.  
E perigoso.  
E, de alguma forma, inevitável.

Larissa já tinha encontrado outras pessoas como eu.  
Volitadores.  
Invisíveis.  
Alguém capaz de copiar o material dos objetos que tocava.

As mesmas que eu encontrara no futuro.

Enquanto Ítalo me distraía, os dois foram atrás delas.

Eu fiquei paralisado no ar.  
Uma mistura de dor e raiva atravessou meu peito.

É impressionante que não importa o que eu faça, o futuro encontra um jeito de convergir para a mesma realidade.  
Talvez fosse esse o plano original dos três desde o começo.  
E Crispim… Crispim era só um pano de fundo.  
Um truque.  
Uma cortina para me distrair.

Talvez seja inevitável que todos saibam sobre as habilidades.  
Talvez tenham esse direito.  
Mas eu vi o futuro.  
Eu vi o caos.  
E eu sei — sei no fundo da alma — que não estão prontos.

Independente do que aconteça, vou até o fim para impedir o que eu vi naquele mundo destruído.

\begin{center}
$\ast$~$\ast$~$\ast$
\end{center}

— Um dirigível bem no meio de Nova Iorque, em plena luz do dia… — murmurou Lúcia enquanto corria, desviando das pessoas. — Como ninguém viu isso chegando?!

Ela era jogada de um lado para o outro pela multidão.  
Quanto mais avançava contra o fluxo, maior era a sensação de que algo gigantesco estava prestes a explodir.

— Não vai pra lá, não! — gritou um senhor grisalho, quase trombando nela. — Você perdeu o juízo?

Ela tentou segurá-lo pelo braço.

— O que está acontecendo? O que está havendo?!

Ele se soltou num impulso, os olhos arregalados.

— Corre, menina! Ainda dá tempo!

Lúcia seguiu.

Carros batidos.  
Gritos ecoando pelas ruas.  
Vidros quebrando.  
E, no meio de tudo isso, uma menina ajoelhada no asfalto, chorando sem conseguir respirar.

— Foi tarde demais… — repetia, como se estivesse presa num ciclo. — Eu não pude fazer nada…

Lúcia se abaixou.

— Ei, ei… olha pra mim. — segurou a menina pelos ombros. — Vem, aqui no meio da rua não é seguro.

A garota parecia atordoada demais para responder.  
Lúcia a ergueu e a levou até um banco.

— Fica aqui. Vou descobrir o que está acontecendo. Eu volto.

— Corre! — gritou um homem de terno azul, fugindo com uma maleta. — Corre, ou você vai morrer aqui!

Ela tentou puxá-lo.

— Terroristas! — gritou uma mulher ao longe. — Eles vão explodir tudo!

E então o chão tremeu.

Uma explosão ecoou sob os pés de todos.

Nova Iorque estava começando a cair.